% Clean-room contest paper baseline (self-written; no upstream template code is vendored).
% Compile with XeLaTeX from the repository root:
%   xelatex -interaction=nonstopmode -output-directory=paper paper/main.tex   (twice)
\documentclass[12pt,a4paper]{ctexart}
\usepackage[margin=2.5cm]{geometry}
\usepackage{amsmath,amssymb,bm}
\usepackage{graphicx}
\usepackage{booktabs,longtable,tabularx,multirow,array}
\usepackage{float}
\usepackage{caption,subcaption}
\usepackage{enumitem}
\usepackage{xcolor}
\usepackage{listings}
\usepackage[hidelinks]{hyperref}   % 正文无彩色：链接与引用均为黑色
\graphicspath{{paper/figures/}}

% ---------------------------------------------------------------------------
% 版面：不要页眉，页码居中放在页脚。
% ctexart 默认用 headings 样式，会把章标题与页码一起放进页眉；
% 这里改用 fancy 并清空页眉，同时把 plain 也重定义一遍（\maketitle 会强制
% 首页用 plain），保证从摘要页开始全篇口径一致。
% ---------------------------------------------------------------------------
\usepackage{fancyhdr}
\pagestyle{fancy}
\fancyhf{}                          % 先清空页眉与页脚
\fancyfoot[C]{\thepage}             % 页码：页脚居中
\renewcommand{\headrulewidth}{0pt}  % 去掉页眉横线
\renewcommand{\footrulewidth}{0pt}  % 去掉页脚横线
\fancypagestyle{plain}{%
  \fancyhf{}%
  \fancyfoot[C]{\thepage}%
  \renewcommand{\headrulewidth}{0pt}%
  \renewcommand{\footrulewidth}{0pt}%
}

% ---------------------------------------------------------------------------
% 圈号（①②…）：xeCJK 默认把 U+2460–24FF 归入西文类，交给 Latin 字体排版，
% 而 Latin 字体没有这两个字形，正文里会缺字或退化成 \textcircled 的西文圈号
% （字重、基线都与中文不齐）。这里把它们改判为 CJK 类，交给正文宋体渲染，
% 于是 ①② 与周围文字同字体、同基线。全篇只有 5.2.2 节用到圈号。
% ---------------------------------------------------------------------------
\xeCJKDeclareCharClass{CJK}{"2460 -> "24FF}

\ctexset{
  section/format=\large\bfseries\heiti,
  section/number=\chinese{section},
  subsection/format=\normalsize\bfseries\heiti,
  subsubsection/format=\normalsize\bfseries\heiti,
}

% ---------------------------------------------------------------------------
% 附录代码清单的排版：整幅细框、行号、按语法着色（灰阶，不喧宾夺主）、
% 自动折行并按缩进对齐，中文注释经 ctex 正常显示。
% ---------------------------------------------------------------------------
\lstdefinestyle{pycode}{
  language=Python,
  basicstyle=\ttfamily\scriptsize,
  keywordstyle=\bfseries\color{black},
  commentstyle=\itshape\color{black!55},
  stringstyle=\color{black!70},
  breaklines=true,
  breakatwhitespace=false,
  postbreak=\mbox{$\hookrightarrow$\space},   % 折行续行标记（不能用 \textcolor，会破坏 listings 的间距计算）
  columns=flexible,
  showstringspaces=false,
  keepspaces=true,
  frame=single,
  framexleftmargin=14pt,
  rulecolor=\color{black!45},
  numbers=left,
  numberstyle=\tiny\color{black!50},
  numbersep=7pt,
  extendedchars=true,
  tabsize=4,
  xleftmargin=20pt,
  xrightmargin=2pt,
  lineskip=1.2pt,
  captionpos=b,
}
\lstset{style=pycode}

% Frozen-number macros (injected by scripts/latex_assembly.py)
__FROZEN_MACROS__

% ---------------------------------------------------------------------------
% 标题块：案例（往届获奖论文）的标题都是一行短句，字号小于 LaTeX 默认的 \LARGE，
% 标题与「摘  要」之间也没有大片空白。这里重定义 \@maketitle 压缩标题占位：
% 标题用三号黑体单行居中，其后只留极小间距，紧接着就是摘要字样。
% 「摘  要」由 \ctexset{abstractname=...} 给出，摘与要之间留一个空白。
% ---------------------------------------------------------------------------
\title{基于径向扩散模型的中药材烘干时长研究}
\author{}
\date{}

\makeatletter
\renewcommand\@maketitle{%
  \vspace*{1.2em}%
  \begin{center}%
    {\zihao{3}\heiti \@title \par}%
  \end{center}%
}
\makeatother

\ctexset{abstractname={摘\hspace{2em}要}}

\linespread{1.4}\selectfont   % 正文行距，对齐获奖论文的 20pt \u884c间距
\captionsetup{font={small},labelsep=quad,justification=centering}
\captionsetup[table]{position=below}
\begin{document}
\maketitle

\begin{abstract}
% 往届获奖论文的摘要一律恰好一页；此处收紧摘要内部的段间距、行距与字号。
% 字体必须在 \begin{abstract} 之内设置：ctex 的 abstract 环境会重置字号。
% 表格 footnotesize 加行距乘子后与正文 12pt 的绝对高度相当，不显突兀。
\begin{small}\linespread{1.22}\selectfont
\setlength{\parskip}{1pt}
烘干是中药材加工中决定成品品质与能耗的关键工序，热风烘干因设备简单、易于控制而普遍采用。
整个过程分为预热平衡与恒温干燥两个阶段，湿润的药材在热风中先升温、再持续失水。
本文以圆柱形中药材在这两个阶段的升温与失水过程为研究对象，分析药材内部温度与水分浓度的
演化规律，并确定烘干所需的时长。四个问题针对同一个药材，
区别只在物性是否随温湿变化、尺寸是否随时间收缩。

\textbf{针对问题一}，建立预热平衡阶段的温湿耦合模型。先由微元法对薄环作热量与水分的守恒分析，
得到一维径向轴对称的导热方程与扩散方程；物性取题目给定的常数，两个方程可以分别求解。
边界取中心的对称条件与表面的第三类对流条件，环境由附件 1 的实测序列分段线性插值驱动。
结果表明该阶段\textbf{以升温为主、失水仅限表层}：$1800$~s 时中心温度已接近环境、
含水率仍为初值，仅表面降至 \qoneCsurfaceoneeightzerozeros；
原因是水分扩散的特征时间远长于本阶段的半小时，水分来不及向内部深处迁移。

\textbf{针对问题二}，改用附录 3 中随温度与含水率变化的物性经验式，建立两者\textbf{双向耦合}的模型：
温度通过扩散系数影响失水，含水率又通过密度、比热容与导热系数影响升温，两个方程必须联立求解。
求解时用三步后向 Euler 启动以抹平预热向恒温切换时初值突变引起的振荡，
其后转为二阶 Crank--Nicolson 格式推进。结果表明干燥由\textbf{内部扩散控制}，
温湿场沿径向单调变化，中心始终滞后于表面。

\textbf{针对问题三}，沿用问题二的模型，以“全域最大干基含水率降至干燥判据以下”为停止条件，
由时间推进直接读出烘干时长。结果表明该判据首次满足的时刻为 \qtwodryingtime，与问题二一致；
判据附近含水率曲线平坦，该时长存在约半分钟量级的不确定度。

\textbf{针对问题四}，药材失水时尺寸收缩，固定网格无法贴合移动的边界。本文改用归一化径向坐标，
把移动的求解区域映射为固定区间，收缩于是表现为一个已知的表观对流项，可在固定网格上求解；
收缩速度由附件 2 数据解析求导得到。为核验移动边界项，本文另设\textbf{解析不变量检验}：
关闭表面传质后，收缩不得改变干基含水率。结果表明\textbf{尺寸收缩不可忽略}：
收缩缩短了扩散路径，其影响压倒了扩散系数的下降，最终烘干时长为 \qfourdryingtimehours，
比不计收缩时更短，故问题四的结论不能由问题三推算。

最后，本文用解析级数解对照、网格与时间加密、独立实现复算与守恒性检验核验结果；
灵敏度分析表明，烘干时长对表面传质系数并不敏感，主要由模型设定决定。
该时长附有成立条件与不确定度，局限在于忽略端面轴向输运与物性经验式未经检验。

\noindent\textbf{关键词}：热风烘干；径向轴对称；有限体积法；移动边界；归一化坐标；参数灵敏度
\end{small}
\end{abstract}

\newpage

% Sections are \input fragments produced by paper-section-writer
__INPUTS__

% AI-use declaration (injected from decision ledgers / ai_use_disclosure.md)
__AI_DECLARATION__

\begin{thebibliography}{99}
__REFERENCES__
\end{thebibliography}

\clearpage
\appendix
\section{支撑材料与程序代码}

\subsection{支撑材料清单}

\begin{table}[H]
  \centering
  \caption{程序文件清单}\label{tab:code}
  \small
  \begin{tabularx}{\textwidth}{@{}llX@{}}
    \toprule
    文件 & 对应问题 & 作用 \\
    \midrule
    \texttt{code/Q1/q1\_common.py} & 问题一 & 公共常数、输入读取、报告点映射与解析级数解 \\
    \texttt{code/Q1/q1\_main.py} & 问题一 & 顶点有限体积离散 + 自适应隐式积分；产出结果表与完整场文件 \\
    \texttt{code/Q1/q1\_baseline.py} & 问题一 & 格心有限体积 + 固定步长 Crank--Nicolson \\
    \texttt{code/Q1/q1\_verifier.py} & 问题一 & 解析级数参照与不变量核验 \\
    \texttt{code/Q2/q2\_main.py} 等 & 问题二 & 含双向耦合的全过程求解、对照实现与检验 \\
    \texttt{code/Q3/q3\_main.py} 等 & 问题三 & 烘干时长的确定、对照实现与检验 \\
    \texttt{code/Q4/q4\_main.py} & 问题四 & 归一化坐标下的移动边界求解，含表观对流项 \\
    \texttt{code/Q4/q4\_baseline.py} & 问题四 & 归一化坐标下的格心有限体积实现 \\
    \texttt{code/Q4/q4\_verifier.py} & 问题四 & 自建级数解与报告值逐格核验 \\
    \texttt{code/figures/} & 全部 & 论文图表的生成脚本（含统一的视觉系统与渲染自检） \\
    \bottomrule
  \end{tabularx}
\end{table}

四个问题的完整计算文件为 \texttt{result1.xlsx}--\texttt{result4.xlsx}，
分别给出题目要求的时间步长与径向间隔下的温度场与水分浓度场。

下文给出核心程序全文。为便于对照正文，每段代码的注释均标出其在论文中的对应位置
（如"论文对应：5.2 节"）。问题二与问题三复用问题一、问题四的求解器，其差异仅在物性
经验式与环境条件，故不再重复列出，完整工程见支撑材料目录 \texttt{code/}。

\subsection{问题一公共模块}
\lstinputlisting[style=pycode, caption={问题一公共模块 \texttt{q1\_common.py}：常数、附件读取、报告点映射与解析级数参照}]
  {code/Q1/q1_common.py}

\subsection{问题一求解程序}
\lstinputlisting[style=pycode, caption={问题一求解程序 \texttt{q1\_main.py}：顶点有限体积离散与自适应 BDF 时间推进}]
  {code/Q1/q1_main.py}

\subsection{问题一检验程序}
\lstinputlisting[style=pycode, caption={问题一检验程序 \texttt{q1\_verifier.py}：Bessel 级数参照与不变量核验}]
  {code/Q1/q1_verifier.py}

\subsection{问题四求解程序}
\lstinputlisting[style=pycode, caption={问题四求解程序 \texttt{q4\_main.py}：\texorpdfstring{$\xi=r/R(t)$}{xi = r/R(t)} 归一化坐标下的移动边界有限体积求解}]
  {code/Q4/q4_main.py}

\subsection{问题四检验程序}
\lstinputlisting[style=pycode, caption={问题四检验程序 \texttt{q4\_verifier.py}：级数参照与交付文件逐格核验}]
  {code/Q4/q4_verifier.py}

\end{document}
